% design.tex — §3 SafeMimic Design
\section{SafeMimic Design}
\label{sec:design}

We now present SafeMimic, a dependency-aware adversarial mimicry framework
that jointly solves the safety and stealthiness problems identified in
\S\ref{sec:introduction}. The key insight is that an operation's
\emph{regularity} (how frequently it appears in normal traces) is neither
necessary nor sufficient for safe camouflage: what matters is whether the
operation's \emph{effect set} intersects with the attack's \emph{dependency
set}. SafeMimic formalizes this insight, constructs a provably safe pool of
candidate operations, and selects from it to minimize behavioral divergence
from the target environment.

% ------------------------------------------------------------------
\subsection{Problem Formulation}
\label{sec:formulation}

\paragraph{Attack chain.}
An attack chain $A = (s_1, s_2, \ldots, s_n)$ is an ordered sequence of
steps, where each step $s_i$ depends on a set of system resources
$\mathrm{deps}(s_i) \subseteq \mathcal{R}$. The resource universe
$\mathcal{R}$ consists of typed resources such as process states, network
connectivity, credentials, file-system objects, and permissions
(\S\ref{sec:dep_extraction}). The \emph{full dependency set} of the attack
is
\begin{equation}
  \Sdep(A) = \bigcup_{i=1}^{n} \mathrm{deps}(s_i).
  \label{eq:sdep}
\end{equation}

\paragraph{Effect set.}
Each candidate benign operation $b$ has an \emph{effect set}
$\Seffect(b) \subseteq \mathcal{R}$: the resources it modifies when
executed. Read-only operations have $\Seffect(b) = \emptyset$.

\paragraph{Goal.}
Given a candidate pool $B_{\text{all}}$ of benign operations and a
baseline behavior profile $\Pnormal$ of the target environment, the
attacker seeks a subset $B \subset B_{\text{all}}$ such that:
\begin{enumerate}
  \item \textbf{Safety:} $\mathrm{success}(A \oplus B) = \mathrm{success}(A)$,
        i.e., the interleaved execution of the attack with the camouflage
        operations preserves attack functionality.
  \item \textbf{Stealthiness:}
        $\DKL(P_{A \oplus B} \| \Pnormal)$ is minimized, where
        $P_{A \oplus B}$ is the behavioral distribution of the
        camouflaged attack trace.
\end{enumerate}

A na\"ive approach would select operations with the highest
\emph{regularity}---the frequency with which they appear in
$\Pnormal$---as ProvNinja~\cite{provninja} does. We show that this
conflates two orthogonal concerns.

\begin{proposition}[Regularity $\neq$ Safety]
\label{prop:regularity}
There exist benign operations $b$ with $\mathrm{freq}(b) \gg 0$ under
$\Pnormal$ yet $\Seffect(b) \cap \Sdep(A) \neq \emptyset$.
\end{proposition}

\begin{proof}
Consider a Linux host where \texttt{kill -9 \textlangle pid\textrangle} is
invoked frequently (e.g., by system services and cron jobs). Its regularity is
high: $\mathrm{freq}(\texttt{kill-process}) \gg 0$. Yet
$\Seffect(\texttt{kill-process}) = \{\textsc{process\_state}\}$, which
intersects $\Sdep(A)$ for any attack chain that depends on a running process
(firefox backdoor, phishing APT, etc.). Executing this operation
during the attack destroys the attacker's foothold and breaks the chain.
\end{proof}

\Cref{tab:safe_pool} (Table~\ref{tab:safe_pool}) summarizes representative
operations with their regularity, effect set, conflict status with a
firefox backdoor attack chain, and membership in $\Bsafe$.

\begin{table}[t]
\centering
\caption{Representative Linux system operations and their safety
analysis against the firefox backdoor attack chain (\S\ref{sec:evaluation}).
Regularity is measured on a benign host trace.
\textbf{High regularity does not imply safety.}}
\label{tab:safe_pool}
\small
\begin{tabular}{@{}lcccc@{}}
\toprule
\textbf{Operation} & \textbf{Reg.} & $\Seffect$ & \textbf{Conflict?} & $\in \Bsafe$? \\
\midrule
\texttt{cat /etc/passwd}    & High & $\emptyset$          & No  & \cmark \\
\texttt{stat /usr/lib/*.so} & High & $\emptyset$          & No  & \cmark \\
\texttt{readdir /proc}      & Med  & $\emptyset$          & No  & \cmark \\
\texttt{cat /var/log/syslog}& High & $\emptyset$          & No  & \cmark \\
\texttt{ls /tmp}            & Med  & $\emptyset$          & No  & \cmark \\
\midrule
\texttt{kill -9 <pid>}      & High & \{process\_state\}   & Yes & \xmark \\
\texttt{iptables -A INPUT}  & Med  & \{network\}          & Yes & \xmark \\
\texttt{systemctl stop}     & Med  & \{service, process\} & Yes & \xmark \\
\texttt{rm -rf /tmp/*}      & Low  & \{file\_system\}     & Yes & \xmark \\
\texttt{touch /tmp/benign}  & Low  & \{file\_system\}     & No  & \cmark \\
\bottomrule
\end{tabular}
\end{table}

% ------------------------------------------------------------------
\subsection{Safe Pool Construction}
\label{sec:safe_pool}

The core of SafeMimic is the construction of a \emph{safe pool}
$\Bsafe$---the subset of candidate operations that provably cannot
interfere with the attack.

\begin{definition}[Safe Pool]
\label{def:safe_pool}
Given an attack chain $A$ with dependency set $\Sdep(A)$ and a candidate
pool $B_{\text{all}}$,
\begin{equation}
  \Bsafe = \bigl\{\, b \in B_{\text{all}} \;\big|\;
           \Seffect(b) \cap \Sdep(A) = \emptyset \,\bigr\}.
  \label{eq:bsafe}
\end{equation}
\end{definition}

\begin{theorem}[Static Safety]
\label{thm:static_safety}
If\/ $b \in \Bsafe$, then $\mathrm{success}(A \oplus b) = \mathrm{success}(A)$.
\end{theorem}

\begin{proof}
Let $b \in \Bsafe$. By \Cref{def:safe_pool},
$\Seffect(b) \cap \Sdep(A) = \emptyset$. Therefore, for every step
$s_i$ in $A$, all resources in $\mathrm{deps}(s_i) \subseteq \Sdep(A)$
remain unmodified after executing $b$, since $b$ only modifies resources
in $\Seffect(b)$. Because each step $s_i$ succeeds if and only if its
dependencies are intact, the entire chain succeeds.
\end{proof}

\begin{corollary}[Composability]
\label{cor:composability}
If\/ $b_1, b_2, \ldots, b_k \in \Bsafe$, then
$\mathrm{success}(A \oplus b_1 \oplus \cdots \oplus b_k) = \mathrm{success}(A)$.
\end{corollary}

\begin{proof}
By induction on $k$. Each $b_j$ satisfies
$\Seffect(b_j) \cap \Sdep(A) = \emptyset$, so the argument of
\Cref{thm:static_safety} applies at each insertion step. No $b_j$
modifies any resource in $\Sdep(A)$, so no $b_j$ affects the
preconditions of any $b_{j'}$ (with $j' > j$) on attack resources.
\end{proof}

\paragraph{Filtering modes.}
SafeMimic supports two filtering granularities:
\begin{itemize}
  \item \textbf{Conservative (type-level):} An operation conflicts if its
        effect contains \emph{any} resource of the same \texttt{ResourceType}
        as an attack dependency, regardless of the specific instance. For
        example, \texttt{kill -9} conflicts with any
        process-state dependency, even if the attack depends on a different process.
  \item \textbf{Precise (instance-level):} Conflict requires an exact
        match on both resource type and instance name. This yields a
        larger $\Bsafe$ but risks false negatives if operations have
        broader effects than their named target.
\end{itemize}
We use conservative mode by default and evaluate the precise mode in
ablation studies (\S\ref{sec:evaluation}).

% ------------------------------------------------------------------
\subsection{Dynamic Safety Space}
\label{sec:dynamic}

The static $\Bsafe$ is conservative: it excludes operations that conflict
with \emph{any} step in the attack, even steps that have already
completed. As the attack progresses through stage~$t$, earlier
dependencies may no longer need protection.

\begin{definition}[Active Dependencies]
\label{def:active_deps}
At attack stage $t$ (after step $s_t$ has executed), the active
dependency set is
\begin{equation}
  \Sdep(A, t) = \bigcup_{i > t} \mathrm{deps}(s_i),
  \label{eq:sdep_t}
\end{equation}
i.e., only the dependencies of steps that remain to be executed.
\end{definition}

\begin{definition}[Dynamic Safe Pool]
\label{def:dynamic_safe_pool}
\begin{equation}
  \Bsafe(t) = \bigl\{\, b \in B_{\text{all}} \;\big|\;
               \Seffect(b) \cap \Sdep(A, t) = \emptyset \,\bigr\}.
  \label{eq:bsafe_t}
\end{equation}
\end{definition}

\begin{theorem}[Dynamic Safety Monotonicity]
\label{thm:dynamic}
For all attack stages $t$:
\begin{enumerate}
  \item $\Bsafe \subseteq \Bsafe(t)$: the dynamic pool is always at
        least as large as the static pool.
  \item If $t_2 > t_1$ and steps $s_{t_1+1}, \ldots, s_{t_2}$ have
        completed between stages $t_1$ and $t_2$, then
        $\Bsafe(t_1) \subseteq \Bsafe(t_2)$.
\end{enumerate}
\end{theorem}

\begin{proof}
(1)~Since $\Sdep(A, t) \subseteq \Sdep(A)$ (we drop completed steps'
dependencies), any $b$ with $\Seffect(b) \cap \Sdep(A) = \emptyset$
also satisfies $\Seffect(b) \cap \Sdep(A, t) = \emptyset$.
(2)~$\Sdep(A, t_2) \subseteq \Sdep(A, t_1)$ because additional steps
have completed between $t_1$ and $t_2$, so
$\Bsafe(t_1) \subseteq \Bsafe(t_2)$ follows by the same argument.
\end{proof}

\paragraph{Implication.}
The monotonicity property means that SafeMimic has \emph{more}
camouflage options at later attack stages, precisely when the attacker is
most active and most in need of cover. In our Linux host experiments
(\S\ref{sec:evaluation}), the dynamic pool grows from 10 operations at
stage~0 to 14 operations by the final stage of a 5-step firefox backdoor
attack.

% ------------------------------------------------------------------
\subsection{Distribution-Matched Insertion}
\label{sec:distribution}

Having constructed a safe pool, SafeMimic must select \emph{which}
operations to insert and \emph{when} to insert them, minimizing the
behavioral divergence from $\Pnormal$ while respecting an edge budget
that prevents edge-count anomaly triggers.

\paragraph{Greedy selection.}
SafeMimic uses a greedy algorithm that iteratively picks the operation
from $\Bsafe$ that most reduces the KL divergence between the
camouflaged trace's operation distribution and $\Pnormal$:
\begin{equation}
  b^{*} = \argmin_{b \in \Bsafe}
           \DKL\!\bigl(P_{\text{selected} \cup \{b\}} \;\|\; \Pnormal\bigr),
  \label{eq:greedy}
\end{equation}
subject to the constraint that the cumulative number of extra provenance
edges does not exceed a budget $E_{\max}$:
\begin{equation}
  \sum_{b \in \text{selected}} \mathrm{edges}(b) \leq E_{\max}.
  \label{eq:budget}
\end{equation}
This budget prevents the camouflage itself from triggering edge-count
anomaly rules (edge count $> \mu + 3\sigma$)~\cite{unicorn,shadewatcher}.

\paragraph{Unique target diversity.}
To avoid inflating the average degree of any single node in the
provenance graph, SafeMimic ensures that each inserted operation targets
a \emph{different} resource. This prevents the creation of
high-fan-out ``hub'' nodes that are easily flagged by structural anomaly
detectors.

\paragraph{Temporal scheduling.}
Rather than inserting all camouflage operations simultaneously (which
creates a detectable burst in time-window-based
detectors~\cite{unicorn}), SafeMimic spreads insertions uniformly across
the attack duration:
\begin{equation}
  t_j = \frac{j \cdot T}{k+1} + \epsilon_j, \qquad
  \epsilon_j \sim \mathcal{N}(0,\, \sigma_{\text{jitter}}^2),
  \label{eq:temporal}
\end{equation}
where $T$ is the total attack duration, $k$ is the number of insertions,
and $\epsilon_j$ is a small Gaussian jitter to avoid perfectly periodic
patterns.

\paragraph{Algorithm overview.}
\Cref{alg:safemimic} presents the complete SafeMimic procedure. The
algorithm takes as input an attack chain, a candidate operation pool, and
a baseline behavior profile, and produces a temporally scheduled
insertion plan that is both provably safe and distribution-matched.

\begin{algorithm}[t]
\caption{SafeMimic: Dependency-Aware Adversarial Mimicry}
\label{alg:safemimic}
\begin{algorithmic}[1]
\Require Attack chain $A = (s_1, \ldots, s_n)$; candidate pool
         $B_{\text{all}}$; baseline profile $\Pnormal$; edge budget
         $E_{\max}$; insertion count $k$
\Ensure Scheduled insertion plan $\mathcal{S}$
\Statex
\State \textbf{// Phase 1: Dependency extraction}
\State $\Sdep(A) \gets \bigcup_{i=1}^{n} \mathrm{deps}(s_i)$
\Statex
\State \textbf{// Phase 2: Safe pool construction}
\State $\Bsafe \gets \{b \in B_{\text{all}} \mid
        \Seffect(b) \cap \Sdep(A) = \emptyset\}$
\Statex
\State \textbf{// Phase 3: Distribution-matched selection}
\State $\text{selected} \gets \emptyset$;~
       $\text{edges\_used} \gets 0$
\For{$j = 1$ \textbf{to} $k$}
  \State $b^{*} \gets \argmin_{b \in \Bsafe \setminus \text{selected}}
          \DKL(P_{\text{selected} \cup \{b\}} \| \Pnormal)$
  \If{$\text{edges\_used} + \mathrm{edges}(b^{*}) > E_{\max}$}
    \State \textbf{break}
  \EndIf
  \State $\text{selected} \gets \text{selected} \cup \{b^{*}\}$
  \State $\text{edges\_used} \mathrel{+}= \mathrm{edges}(b^{*})$
\EndFor
\Statex
\State \textbf{// Phase 4: Temporal scheduling}
\State $T \gets \text{attack duration}$
\For{$j = 1$ \textbf{to} $|\text{selected}|$}
  \State $t_j \gets \frac{j \cdot T}{|\text{selected}| + 1}
          + \epsilon_j$,~~$\epsilon_j \sim \mathcal{N}(0, \sigma^2)$
\EndFor
\State $\mathcal{S} \gets \{(t_j,\, b_j) \mid b_j \in \text{selected}\}$
\State \Return $\mathcal{S}$
\end{algorithmic}
\end{algorithm}

\paragraph{Complexity.}
The greedy loop (\Cref{alg:safemimic}, lines 7--13) runs $k$ iterations,
each evaluating $|\Bsafe|$ candidates. Computing $\DKL$ over the
discrete operation-type distribution takes $O(|\mathcal{T}|)$ time, where
$\mathcal{T}$ is the set of operation types. The overall complexity is
$O(k \cdot |\Bsafe| \cdot |\mathcal{T}|)$, which is negligible compared
to the attack execution itself (typically under 100\,ms for our
17-operation pool).
